\documentclass[11pt,a4paper]{article}
\usepackage{fontspec,unicode-math,graphicx,xcolor,colortbl,array,tabularx,longtable,geometry,fancyhdr,enumitem,hyperref}
\geometry{left=20mm,right=20mm,top=19mm,bottom=20mm}
\setmainfont{BeVietnamPro_Regular.ttf}[Path=../assets/,BoldFont=BeVietnamPro_SemiBold.ttf]
\setmathfont{STIXTwoMath.ttf}[Path=../assets/]
\definecolor{ink}{HTML}{3E3A35}\definecolor{muted}{HTML}{766F66}\definecolor{grid}{HTML}{DFD7CA}\definecolor{paper}{HTML}{FBFAF8}\definecolor{head}{HTML}{F8F5F0}\definecolor{red}{HTML}{BF4240}\definecolor{brandred}{HTML}{EF5350}\definecolor{yellow}{HTML}{FFCA28}\definecolor{focus}{HTML}{FFF9E9}
\pagecolor{paper}\color{ink}\arrayrulecolor{grid}\setlength\arrayrulewidth{.4pt}\setlength\tabcolsep{7pt}\renewcommand\arraystretch{1.5}\setlength\parindent{0pt}\setlength\parskip{8pt}\linespread{1.12}
\pagestyle{fancy}\fancyhf{}\fancyhead[L]{\scriptsize\color{muted}ZO MATH / HỆ THIẾT KẾ 0.4}\fancyhead[R]{\scriptsize\color{muted}06.10.2026}\fancyfoot[L]{\scriptsize\color{muted}Nhận ra ZO. Tập trung vào Toán.}\fancyfoot[R]{\scriptsize\thepage}\renewcommand{\headrulewidth}{.3pt}\renewcommand{\footrulewidth}{0pt}
\hypersetup{colorlinks=true,urlcolor=red,linkcolor=red}\setlist{nosep,leftmargin=1.5em}
\newcommand{\tagline}[1]{{\small\color{muted}\MakeUppercase{#1}}\par}
\newcommand{\titlepagehead}[2]{\tagline{#1}{\fontsize{27}{32}\selectfont\bfseries #2}\par\vspace{6mm}}
\newcommand{\why}[1]{\par{\small\color{muted}\textbf{Vì sao.} #1}\par}
\newcommand{\pic}[2]{\begin{center}\includegraphics[width=#2\linewidth]{#1}\end{center}}
\begin{document}
\tagline{Bảng mẫu / 0.4}\section*{05 · Bảng dấu}
\begin{tabularx}{\linewidth}{|>{\raggedright\arraybackslash}p{.28\linewidth}|>{\centering\arraybackslash}X|>{\centering\arraybackslash}X|>{\centering\arraybackslash}X|>{\centering\arraybackslash}X|>{\centering\arraybackslash}X|}\hline
\rowcolor{head}$x$ & $(-\infty;-1)$ & $-1$ & $(-1;1)$ & $1$ & $(1;+\infty)$\\\hline
$x+1$ & $-$ & $0$ & $+$ & $+$ & $+$\\\hline
$x-1$ & $-$ & $-$ & $-$ & $0$ & $+$\\\hline
\rowcolor{focus}$p^{\prime}(x)=3(x+1)(x-1)$ & $+$ & $0$ & $-$ & $0$ & $+$\\\hline
\end{tabularx}\why{Giữ ô riêng cho khoảng và mốc. Hàng kết quả dùng yellow-1; dấu cộng/trừ vẫn viết bằng LaTeX và cùng màu chữ, không suy nghĩa từ màu.}\end{document}